\section{Forecasting Method}

The modeling design follows the EDA: Revenue and COGS are related but not identical business processes. Revenue is a demand process; COGS is a margin process.

\paragraph{Revenue model.}
Revenue is decomposed into a period-level demand anchor and a daily seasonal allocation:
\[
  \widehat{\text{Revenue}}_t =
  \widehat{\text{Demand level}}_{\text{period}(t)}
  \times
  \widehat{\text{Seasonal share}}_t.
\]
The demand anchor uses historical Revenue, calendar recency, and train-derived regime priors. The daily allocation uses month, day-of-week, month-day, boundary-day, and historical seasonal profiles. H1 receives more shape confidence because its historical pattern is more stable; H2 is regularized more strongly.

\paragraph{COGS ratio head.}
Instead of directly scaling COGS from Revenue, the model estimates a COGS/Revenue ratio prior from historical month, discount, promotion, and regime evidence:
\[
  \widehat{\text{COGS}}_t =
  \widehat{\text{Revenue}}_t
  \times
  \widehat{r}_t,\quad
  r_t = \text{COGS}_t / \text{Revenue}_t.
\]
This separates margin pressure from demand level and makes the model robust to the fact that promotion-heavy days can have normal or high Revenue but unusually compressed gross margin.

\paragraph{Daily allocation and stabilization.}
The final configuration uses a conservative month-preserving daily stabilizer for H2 COGS ratio risk. It adjusts daily COGS allocation toward recent even-year historical ratio patterns while preserving period/month totals. This is intended to reduce tail-day error and avoid overreacting to unstable H2 daily spikes.

\paragraph{Feature policy.}
Features are grouped by availability:
\begin{itemize}
  \item \textbf{Known future:} calendar, month, day-of-week, month boundaries, and algorithmic lunar-calendar positions derived from \texttt{Date}.
  \item \textbf{Historical priors:} recurring promotion intensity, historical ratio profiles, demand-regime priors, source/category/region mix priors.
  \item \textbf{Excluded as realized future inputs:} future orders, realized traffic, future returns, future inventory states, and any unavailable post-2022 operational measurements.
\end{itemize}

\paragraph{Explainability.}
The key model drivers are interpretable by construction: demand regime, calendar seasonality, recurring promotion pressure, and COGS/Revenue ratio. For tree-based anchor components, feature importance and partial dependence checks are used as equivalent explanations to verify that the model relies on calendar, recency, and train-derived promotion/ratio signals rather than unavailable future information.
